\section{Representation Learning and Classification Baselines}

\subsection{Feature Backbone, Projection Head, and Classification Baseline}
After the test set is frozen, recognition and representation learning are developed on a shared backbone and nonlinear projection head. For downstream tasks using metric learning or self-supervised learning, a ConvNeXt-Tiny network~\cite{convnext} serves as the feature backbone, producing a feature vector $u_i$. This vector is passed through a two-layer fully connected projection head with ReLU activations to produce an $L_2$-normalized embedding $v_i\in\mathbb{R}^{256}$:
\begin{equation}
\tilde{v}_i=W_2\,\mathrm{ReLU}(W_1u_i+b_1)+b_2,\qquad
v_i=\frac{\tilde{v}_i}{\|\tilde{v}_i\|_2}.
\label{eq:projection_head}
\end{equation}

As a classification reference, we train an Image Classification Baseline using ConvNeXt-Tiny with a direct fully connected classifier head predicting probabilities over 18 wood species. To address class imbalance, the model is trained with Focal Loss rather than standard cross-entropy:
\begin{equation}
L_{\text{Focal}} = -\frac{1}{B} \sum_{i=1}^{B} \alpha (1 - p_{t,i})^\gamma \log(p_{t,i}),
\label{eq:focal_loss}
\end{equation}
where $p_{t,i}$ is the predicted probability for the ground-truth class of sample $i$, $\gamma$ is the focusing parameter that downweights easy samples, and $\alpha$ is a balancing factor. In all experiments, $\gamma = 2.0$ and $\alpha = 0.25$. Approximately 90\% of the backbone layers are frozen to limit overfitting on the limited-specimen dataset. The optimizer is AdamW with an initial learning rate of $5\times10^{-4}$ and weight decay $10^{-2}$, combined with a Cosine Annealing scheduler. All metric learning and self-supervised learning methods share the same manifest, split, augmentation pipeline, and validation selection rule; observed differences between methods therefore cannot be attributed to split variation.

\subsection{Deep Metric Learning Loss Functions}
To evaluate representation learning approaches, we benchmark four core distance-based loss functions under a unified training setup with the same backbone, projection head, augmentation, and CEGS-Split. Shared notation: $\mathbf{z}_i\in\mathbb{R}^{D}$ is the $L_2$-normalized embedding, $y_i$ is the species label, $d_{ij}=\|\mathbf{z}_i-\mathbf{z}_j\|_2$ is the Euclidean distance, $S_{ij}=\mathbf{z}_i^T\mathbf{z}_j$ is cosine similarity, and $[u]_+=\max(0,u)$.

\subsubsection{Standard Contrastive Loss}
Contrastive Loss optimizes pairwise distances directly:
\begin{equation}
L_{\text{contrastive}}=(1-q_{ij})\tfrac{1}{2}d_{ij}^{2}+q_{ij}\tfrac{1}{2}[m-d_{ij}]_+^2,
\end{equation}
where $q_{ij}=1$ if the two samples belong to the same species and $q_{ij}=0$ otherwise.

\subsubsection{Standard Triplet Loss}
Triplet Loss optimizes the relation among an anchor $a$, a positive $p$ from the same species, and a negative $n$ from a different species, enforcing a margin $\alpha$ between the positive and negative distances:
\begin{equation}
L_{\text{triplet}}=\frac{1}{|\mathcal{T}|}\sum_{(a,p,n)\in\mathcal{T}}\left[d(a,p)^2-d(a,n)^2+\alpha\right]_+.
\end{equation}

\subsubsection{Semi-hard Triplet Loss}
Semi-hard Triplet Loss selects negatives that are neither trivial nor overly noisy. For an anchor--positive pair $(a,p)$, the negative $n$ is selected when:
\begin{equation}
d(a,p)^2<d(a,n)^2<d(a,p)^2+\alpha.
\end{equation}
The loss retains the standard triplet form over the semi-hard set $\mathcal{T}_{sh}$:
\begin{equation}
L_{\text{semi-hard}}=\frac{1}{|\mathcal{T}_{sh}|}\sum_{(a,p,n)\in\mathcal{T}_{sh}}\left[d(a,p)^2-d(a,n)^2+\alpha\right]_+.
\end{equation}

\subsubsection{Angular Loss}
Angular Loss optimizes the angular relationship within triplets rather than Euclidean distances alone:
\begin{equation}
L_{\text{angular}}=\frac{1}{|\mathcal{T}|}\sum_{(a,p,n)\in\mathcal{T}}\log\left[1+e^{\Delta_{\theta}(a,p,n)}\right],
\end{equation}
where:
\begin{equation}
\begin{aligned}
\Delta_{\theta}(a,p,n) = &4\tan^2\!\alpha\,(\mathbf{z}_a+\mathbf{z}_p)^T\mathbf{z}_n \\
&- 2(1+\tan^2\!\alpha)\,\mathbf{z}_a^T\mathbf{z}_p.
\end{aligned}
\end{equation}

\subsection{Self-Supervised Learning Loss Functions}
In self-supervised learning (SSL), the model optimizes representations through randomly augmented views of each image without using class labels. We evaluate four SSL methods.

\subsubsection{SimCLR}
SimCLR~\cite{simclr} minimizes the InfoNCE contrastive loss over augmented pairs. For a batch of $N$ images, each augmented into two views, the loss for positive pair $(i,j)$ is:
\begin{equation}
\ell(i, j) = -\log \frac{e^{\mathbf{z}_i^T\mathbf{z}_j/\tau}}{\sum_{k=1}^{2N} \mathbb{I}_{[k \ne i]}\, e^{\mathbf{z}_i^T\mathbf{z}_k/\tau}},
\end{equation}
where $\tau$ is the temperature hyperparameter. The overall loss averages over all positive pairs:
\begin{equation}
L_{\text{SimCLR}} = \frac{1}{2N} \sum_{k=1}^{N} \left[ \ell(2k-1, 2k) + \ell(2k, 2k-1) \right].
\end{equation}

\subsubsection{BYOL}
BYOL~\cite{byol} uses an online network with parameters $\theta$ and a target network with parameters $\xi$. The online network predicts the target network's representation of a different augmented view. The loss is the mean squared error between normalized predictions:
\begin{equation}
L_{\text{BYOL}} = \left\| \bar{q}_\theta - \bar{z}'_\xi \right\|_2^2 = 2 - 2\,\bar{q}_\theta^T \bar{z}'_\xi,
\end{equation}
where $\bar{q}_\theta$ and $\bar{z}'_\xi$ are $L_2$-normalized. Collapse is prevented by applying a stop-gradient to the target network, whose parameters are updated via exponential moving average:
\begin{equation}
\xi \leftarrow m\,\xi + (1-m)\,\theta,\quad m\in[0,1].
\end{equation}

\subsubsection{SimSiam}
SimSiam~\cite{simsiam} simplifies BYOL by removing the momentum target network (setting $\xi=\theta$) and relying solely on a stop-gradient operator and an asymmetric predictor to prevent collapse. The symmetric loss is:
\begin{equation}
L_{\text{SimSiam}} = \tfrac{1}{2}\,\mathcal{D}(p_1, \mathrm{sg}(z_2)) + \tfrac{1}{2}\,\mathcal{D}(p_2, \mathrm{sg}(z_1)),
\end{equation}
where $p_k=q_\theta(z_k)$ is the predictor output, $\mathrm{sg}$ denotes stop-gradient, and $\mathcal{D}(p,z)=-p^Tz/(\|p\|_2\|z\|_2)$.

\subsubsection{Barlow Twins}
Barlow Twins~\cite{barlow} applies a redundancy-reduction objective to the cross-correlation matrix $C$ of size $D\times D$, computed along the batch dimension between projections $\mathbf{z}^A$ and $\mathbf{z}^B$:
\begin{equation}
C_{ij} = \frac{\sum_{b=1}^{B} z_{b,i}^A z_{b,j}^B}{\sqrt{\sum_{b=1}^{B} (z_{b,i}^A)^2}\,\sqrt{\sum_{b=1}^{B} (z_{b,j}^B)^2}}.
\end{equation}
The loss drives diagonal entries toward 1 (invariance term) and off-diagonal entries toward 0 (redundancy-reduction term):
\begin{equation}
L_{\text{BT}} = \sum_{i}(1 - C_{ii})^2 + \lambda \sum_{i}\sum_{j\neq i} C_{ij}^2,
\end{equation}
where $\lambda>0$ controls the relative weight of the redundancy term.
